% Pursuing Stacks (cotes 134-2 to 134-8): archive cote 134-4, site page 31 (PDF page 32), typescript page 119.
% typescript. Draft compared: ps3.tex:1532-1597.
% AI-transcribed by Claude Opus 5.5 (claude-opus-5-5, adaptive thinking, effort high), Message Batches API, batch ps2 (msgbatch_01NVwr1Yuo8qJf5rykeGBXPK), 2026-11-12; not yet checked by eye.
\scanpage{134-4/31}
\mspage{119}

by the subset $\hand{\Sigma}_{hc}$ of it's contractible elements, which is just another
way of saying that this $\hand{\Sigma}_h \hand{\in} \mathrm{Hom}_3(M)$ comes already from a contractibility
structure $C_h \hand{\in} \mathrm{Hom}_4(M)$ ! Thus definitely the uniformity of \inserted{formal} relationships
between successively occuring notions seems broken here, namely there
does not seem to be any natural retraction $r_{43}$ of $\mathrm{Hom}_3(M)$ unto the subset
$\mathrm{Hom}_4(M)$ . For the least, if there is such a retraction, it's definition
should be presumably a somewhat more delicate one than the first that \xout{xxxx}
comes to mind. I will not pursue this matter any further now, as it is
not clear if we'll need it later.

\deleted{In case M admits finite products,} \over{i}{I}t is clear that \xout{xxx} for any
weak homotopy interval structure $\hand{\Sigma}_h$ on M, $\hand{\Sigma}_h$ is stable under the natu-%
ral notion of \inserted{finite} products\xout{x} of intervals (in the sense of the category structu-%
te \xout{\unsure{of}} \uline{Int}(M)). \over{w}{W}e saw already that this is handy, as the consideration
of products of intervals allow\hand{s} to show that the family of homotopy
relations $\underset{\underline{I}}{\hand{\sim}}$ in Fl(M), for variable $\underline{I}$ in $\hand{\Sigma}_h$ , is ``filtrant décrois-%
sant'' , so we get the relation $\underset{h}{\hand{\sim}}$ as the filtering direct limit or
union of the more elementary relations $\underset{\underline{I}}{\hand{\sim}}$ . Similarly, if $C_h$ \inserted[typed]{$\hand{\subset}$ Ob(M)} is a
contractibility structure on M, $C_h$ is stable under finite products.

\uline{Remark} From the way we've been working so far with homotopy notions, it
would seem that we're only interested \inserted{here} in homotopy notions which stem
from a structure in $\mathrm{Hom}_3(M)$, namely which can be described in terms of a
notion of weak homotopy intervals. The focus on contractibility has set
in only lately, and it is too soon to be sure whether we'll be working
only with homotopy structures on M which can be described in terms of a
contractibility notion, namely which are in $\mathrm{Hom}_4(M)$. In the cases I've
had in mind so far, it turns out \xout{xxxx} it seems that the homotopy notions
dealt with do come from a structure in $\mathrm{Hom}_4(M)$, \inserted{i.e. from} a contractibility struc-%
ture.

E) \uline{Generating sets of \xout{inter} weak homotopy intervals}. \inserted{\uline{Contractors}.} Let $\hand{\Sigma}_h$ be a
weak homotopy interval structure on M. \inserted{(This implies we assume M stable under finite products)} A subset $\hand{\Sigma}^{\circ}_{h}$ is called \uline{generat} -%
\uline{ting} , if $\hand{\Sigma}_{\unsure{h|H}}$ is just it's saturation (cf. C) above) i.e. for any $\underline{I}$ in
$\hand{\Sigma}_h$ , the two endpoints can be joined by a chain\unsure{\xout{,}} as in (Sat 3), invol-%
ving only intervals in $\hand{\Sigma}^{\circ}_{h}$ . This implies that all homotopy notions
dealt with \xout{xxx} so far can be checked directly in terms of intervals in
$\hand{\Sigma}^{\circ}_{h}$ . We've met a particular case of this before, when $\hand{\Sigma}^{\circ}_{h}$ is reduced
to just one element $\underline{I}$ - we then called $\underline{I}$ ``characteristic'' , but
``\uline{generating} \inserted[typed]{\uline{weak homotopy interval}''} now would seem the more appropriate expression. Even when

\margin{31}
